%source mathadata
\begin{multicols}{2}
		Une droite $d$ est donnée par sa forme réduite : $$\displaystyle d : y=\frac{2}{3}x-\frac{7}{3}.$$ Les points au-dessus de la droite sont considérés comme verts ;  et en dessous ou sur la droite,  comme bleus.
		\begin{MADenumerate}
			\item Pour chaque point, déterminer si le point appartient à la droite $d$ ou non par un calcul, sans regarder le graphique.

			(a) $A(5;1)$ \; (b) $B(3;2)$ \; (c) $C(5;-3)$

			(d) $D(2;-1)$ \; (e) $E(0,5;-2)$ \; (f) $F(-1; -2)$

			\item Tracer la droite $d$ et les points donnés afin de vérifier vos résultats.

			\item Déterminer les couleurs des 6 points.

			\columnbreak

			\item (Bonus) Proposer un point $G$ à coordonnées entières qui appartient à $d$ et un point $H$ d'ordonnée comprise entre -0,5 et 0,5 qui appartient à $d$. Quelles sont leurs couleurs ?

		\end{MADenumerate}

		%\begin{center}
		\begin{figure}[H]\centering
			\definecolor{wrwrwr}{rgb}{0.380,0.380,0.380}
			\definecolor{rvwvcq}{rgb}{0.082,0.396,0.753}
			\begin{tikzpicture}[line cap=round,line join=round,>=triangle 45,x=1cm,y=1cm,scale=0.7]
				\begin{axis}[
						x=1cm,y=1cm,
						axis lines=middle,
						ymajorgrids=true,
						xmajorgrids=true,
						xmin=-1.5,
						xmax=7.5,
						ymin=-4.5,
						ymax=4.5,
						xtick={-1,...,7},
						ytick={-4,...,4},
						tick label style={font=\Large},
					]

					%\draw[domain=-1:7] plot(\x,{(2*\x-7)/3});

					%\node[above,font=\LARGE] at (axis cs:6,2) {d};

					% Annotations pour les zones de classification
					\node[font=\normalsize, text=green!60!black] at (axis cs:3,3.5) {Au-dessus : verts};
					\node[font=\normalsize, text=blue!60!black] at (axis cs:3,-3.5) {En dessous : bleus};

				\end{axis}
			\end{tikzpicture}
		\end{figure}
		%\end{center}

	\end{multicols}
